\chapter{Uogólnione twierdzenia Riccatiego}
\label{ch:lin26-general-riccati}
\index{równanie Riccatiego!stabilizowalność}
\index{wykrywalność!równanie Riccatiego}
\index{sterowanie optymalne!koszt półokreślony}

\noindent\textbf{Cel rozdziału.}
W Rozdziale~\ref{ch:lin24-riccati}
udowodniliśmy twierdzenia Riccatiego
przy stabilności macierzy \(A\)
i dodatniej określoności \(Q\).
W wielu układach macierz otwartej
pętli jest jednak niestabilna,
a koszt nie karze każdej współrzędnej
stanu. Wykażemy, że wystarczają
\emph{stabilizowalność} pary \((A,B)\)
i \emph{wykrywalność} pary
obserwacyjnej, które zostały
zdefiniowane i scharakteryzowane
w Rozdziale~\ref{ch:lin25-control}.
Przedstawimy pełny dowód
istnienia i jednoznaczności
rozwiązania stabilizującego,
wzór na wartość optymalną,
przykłady z niestabilną otwartą
pętlą i kontrprzykłady
wyjaśniające rolę założeń.

Niech \(n,m\ge1\).
W całym rozdziale
\(A\in\operatorname{Mat}_{n\times n}(\mathbb C)\),
\(B\in\operatorname{Mat}_{n\times m}(\mathbb C)\),
\(Q=Q^*\ge0\) i
\(R=R^*>0\) są macierzami
odpowiednich rozmiarów \(n\times n\)
i \(m\times m\).
Przez \(P\ge0\)
oznaczamy nieujemną określoność
hermitowskiej macierzy \(P\).
Symbol \(\sigma(T)\) oznacza
zbiór wartości własnych \(T\),
a \(\rho(T)\) jego promień spektralny.
Przez \(u=-Kx\)
oznaczamy sprzężenie zwrotne,
które wprowadza macierz
zamkniętej pętli \(A-BK\).

\section{Założenia i koszt półokreślony}
\label{sec:lin26-hypotheses}

\begin{lemat}[Rozkład macierzy kosztu]
\label{lem:lin26-factor}
Dla każdej hermitowskiej
\(Q\in\operatorname{Mat}_{n\times n}(\mathbb C)\),
\(Q\ge0\), istnieje
hermitowska macierz \(C\ge0\)
taka, że \(C^*C=Q\).
Ponadto dla każdego
\(x\in\mathbb C^n\) mamy
\[
 x^*Qx=\|Cx\|_2^2,\qquad
 \ker C=\ker Q.
\]
Jeżeli \(C_1^*C_1=C_2^*C_2=Q\)
dla macierzy o dowolnej liczbie
wierszy, to pary
\((C_1,A)\) i \((C_2,A)\)
są wykrywalne w danym rodzaju czasu
jednocześnie albo żadna z nich
nie jest wykrywalna.
\end{lemat}
\begin{proof}
Z twierdzenia spektralnego
z Rozdziału~\ref{ch:lin14-spektralne}
istnieje unitarna \(U\)
i liczby \(q_j\ge0\)
takie, że
\(Q=U\operatorname{diag}(q_1,\ldots,q_n)U^*\).
Połóżmy
\(C=U\operatorname{diag}(\sqrt{q_1},
\ldots,\sqrt{q_n})U^*\).
Jest to macierz hermitowska,
nieujemnie określona, a
\(C^*C=C^2=Q\).
Równość
\(x^*Qx=(Cx)^*(Cx)\)
daje wzór na koszt.

Jeśli \(Cx=0\), to
\(Qx=C^*Cx=0\).
Jeśli \(Qx=0\),
to \(\|Cx\|_2^2=x^*Qx=0\),
więc \(Cx=0\).
To dowodzi równości jąder.
Analogicznie
\(\ker C_1=\ker Q=\ker C_2\).
Na mocy kryterium PBH
z Twierdzenia~\ref{thm:lin25-detectability}
wykrywalność zależy tylko od
tego, czy dla każdej niestabilnej
wartości własnej \(\lambda\)
istnieje niezerowy
\(v\in\ker(\lambda I-A)\cap\ker C_i\).
Równość jąder oznacza,
że oba warunki są identyczne.
\end{proof}

\begin{definicja}[Warunki dla ogólnego LQR]
\label{def:lin26-assumptions}
Niech \(C\) będzie dowolną
macierzą spełniającą
\(C^*C=Q\).
Przyjmujemy dwa warunki:
\begin{enumerate}[label=\textup{(\roman*)}]
 \item para \((A,B)\) jest
 \emph{stabilizowalna}
 w rozważanym rodzaju czasu;
 \item para \((C,A)\) jest
 \emph{wykrywalna}
 w tym samym rodzaju czasu.
\end{enumerate}
W pierwszym przypadku
istnieje \(K_0\)
stabilizujące \(A-BK_0\).
W drugim niestabilne
wartości własne \(A\)
nie mają niezerowych
wektorów własnych w \(\ker Q\).
Z
Lematu~\ref{lem:lin26-factor}
warunek (ii) nie zależy
od wyboru rozkładu
\(Q=C^*C\).
\end{definicja}

\begin{uwaga}[Dlaczego nie wymagamy \(Q>0\)]
\label{rem:lin26-weak-q}
Jeśli \(Q>0\), to \(C\)
jest odwracalna, zatem
para \((C,A)\) jest obserwowalna,
a w szczególności wykrywalna.
Warunek wykrywalności
jest więc istotnie słabszy.
Pozwala on nie karać
stabilnych kierunków, których
obserwacja nie rozróżnia.
Nie będziemy jednak wnioskować,
że rozwiązanie Riccatiego
musi być dodatnio określone:
w ogólnym przypadku otrzymamy
tylko \(P\ge0\).
\end{uwaga}

\section{Dwa oszacowania stabilnego sprzężenia}
\label{sec:lin26-feedback-bound}

\begin{lemat}[Koszt stabilizującego sterowania]
\label{lem:lin26-feedback-bound}
Załóżmy, że para \((A,B)\)
jest stabilizowalna.
Wybierzmy macierz \(K_0\)
taką, aby \(D=A-BK_0\)
była stabilna w odpowiednim
rodzaju czasu i połóżmy
\(H=Q+K_0^*RK_0\ge0\).

W czasie dyskretnym szereg
\[
 W_{\mathrm d}
 =\sum_{k=0}^{\infty}(D^*)^kHD^k
\]
jest bezwzględnie zbieżny,
daje macierz hermitowską
\(W_{\mathrm d}\ge0\)
i spełnia
\[
 W_{\mathrm d}-D^*W_{\mathrm d}D=H.
\]
Dla trajektorii
\(x_{k+1}=Dx_k\)
koszt sterowania \(u_k=-K_0x_k\)
na nieskończonym horyzoncie
wynosi \(x_0^*W_{\mathrm d}x_0\).

W czasie ciągłym całka
\[
 W_{\mathrm c}
 =\int_0^\infty e^{tD^*}He^{tD}\,\mathrm dt
\]
jest bezwzględnie zbieżna,
daje \(W_{\mathrm c}=W_{\mathrm c}^*\ge0\)
i spełnia
\[
 D^*W_{\mathrm c}+W_{\mathrm c}D=-H.
\]
Dla trajektorii \(x'=Dx\)
koszt sterowania \(u=-K_0x\)
jest równy \(x_0^*W_{\mathrm c}x_0\).
\end{lemat}
\begin{proof}
W czasie dyskretnym
\(\rho(D)<1\), więc
z Lematu~\ref{lem:lin23-powers}
mamy \(\|D^k\|\le Mr^k\)
dla pewnych \(M>0\), \(0<r<1\).
Wobec tego
\[
 \|(D^*)^kHD^k\|
 \le M^2\|H\|r^{2k},
\]
co zapewnia zbieżność.
Każdy składnik jest
hermitowski i nieujemnie
określony, dlatego
\(W_{\mathrm d}\ge0\).
Odjęcie od szeregu
tego samego szeregu
z indeksem przesuniętym
daje równanie dla \(W_{\mathrm d}\).
Ponieważ
\(u_k^*Ru_k+x_k^*Qx_k=x_k^*Hx_k\),
sumowanie kosztów
daje wskazaną formę kwadratową.

W czasie ciągłym
wszystkie wartości własne \(D\)
mają ujemne części rzeczywiste.
Z postaci Jordana lub
ciągłego twierdzenia Lapunowa
z Rozdziału~\ref{ch:lin22-sylvester}
wynika oszacowanie
\(\|e^{tD}\|\le Me^{-\alpha t}\)
dla \(t\ge0\)
i pewnych \(M,\alpha>0\).
Zatem
\(\|e^{tD^*}He^{tD}\|
 \le M^2\|H\|e^{-2\alpha t}\),
więc całka jest zbieżna
i definiuje macierz
nieujemnie określoną.
Różniczkując macierz
\(e^{tD^*}He^{tD}\)
i całkując jej pochodną,
otrzymujemy
\[
 \begin{aligned}
 D^*W_{\mathrm c}+W_{\mathrm c}D
 &=\int_0^\infty
     \frac{\mathrm d}{\mathrm dt}
       (e^{tD^*}He^{tD})\,\mathrm dt\\
 &=\lim_{t\to\infty}
       e^{tD^*}He^{tD}-H=-H.
 \end{aligned}
\]
Koszt ma postać
\(\int_0^\infty x_0^*e^{tD^*}
He^{tD}x_0\,\mathrm dt\);
przesunięcie stałych
wektorów poza całkę
daje \(x_0^*W_{\mathrm c}x_0\).
\end{proof}

\begin{uwaga}[Granica dodatnia nie musi być ściśle dodatnia]
\label{rem:lin26-feedback-semi}
W poprzednim lemacie
\(H\) może mieć niezerowe
jądro. Nie można więc
stosować twierdzeń Lapunowa
w wersji wymagającej \(H>0\)
do wniosku o stabilności:
stabilność macierzy \(D\)
była tu \emph{założeniem},
a nie wynikiem.
Dopiero wykrywalność
umożliwi wyprowadzenie
stabilności nowego sprzężenia
z równania Riccatiego
o prawostronnej macierzy
jedynie nieujemnie określonej.
\end{uwaga}
